% Relatório técnico do trabalho prático de Processamento Digital de Imagens.
% Compilar com:  cd report && pdflatex relatorio.tex && pdflatex relatorio.tex   (ou ./compilar.sh)
% As figuras (figuras/) e as tabelas (gerado/) vêm de `python report/gerar_figuras.py`.
\documentclass[12pt,a4paper]{article}

\usepackage[utf8]{inputenc}
\usepackage[T1]{fontenc}
\usepackage{lmodern}
\usepackage[brazilian]{babel}
\usepackage[a4paper,margin=2.5cm]{geometry}
\usepackage{graphicx}
\usepackage{booktabs}
\usepackage{tabularx}
\usepackage{caption}
\usepackage{float}
\usepackage{amsmath,amssymb}
\usepackage{xcolor}
\usepackage{enumitem}
\usepackage{setspace}
\usepackage{microtype}
\usepackage{tikz}
\usetikzlibrary{positioning,shapes.geometric,arrows,calc}
\usepackage[hidelinks]{hyperref}

\setstretch{1.15}
\setlength{\parindent}{1.25cm}
\setlength{\parskip}{0.2em}
\captionsetup{font=small,labelfont=bf}
\setlist{itemsep=0.15em,topsep=0.3em}
\renewcommand{\arraystretch}{1.12}

\newcommand{\Sone}{S\textsubscript{1}}
\newcommand{\Stwo}{S\textsubscript{2}}

\hypersetup{pdftitle={Reconstrução de sinais de áudio cardíacos a partir de traçados em imagem},
            pdfauthor={Alex Tavares Cordeiro}}

\begin{document}

\begin{titlepage}
\centering
{\large Universidade Federal da Paraíba\par}
{\large Centro de Informática\par}
\vspace{0.4cm}
{\large Introdução ao Processamento Digital de Imagens --- Período 2026.2\par}
{\large Professor: Leonardo Vidal\par}

\vspace{3.2cm}
{\LARGE\bfseries Reconstrução automática de sinais de áudio cardíacos a partir de traçados em imagem\par}
\vspace{0.8cm}
{\large Trabalho prático\par}

\vspace{3.5cm}
{\large Aluno: Alex Tavares Cordeiro\par}

\vfill
{\large Novembro de 2026\par}
\end{titlepage}

\begin{abstract}
\noindent Este trabalho reconstrói, a partir de uma imagem, três sinais de áudio de sons cardíacos: um sem sopro (Normal), um com o sopro inocente de Still e um com o sopro patológico de comunicação interventricular (VSD). A imagem é a Figura~1 do artigo de Tomas et al.\ (2011), que mostra os três sinais como gráficos em função do tempo. O programa, escrito em Python, isola cada gráfico, remove a grade, os eixos, os textos e as anotações, extrai a curva, calibra o eixo do tempo e o da amplitude, e grava os arquivos \texttt{normal.wav}, \texttt{still.wav} e \texttt{vsd.wav} (5000 amostras por segundo, 16 bits, mono). A filtragem espacial usa uma função de correlação bidimensional escrita do zero, no domínio do espaço, com preenchimento por zeros. A remoção da grade usa uma medida de cor que não depende do fundo, e a extração usa uma regra por coluna que preserva a amplitude dos picos. A análise dos sinais reconstruídos reproduz sete de oito afirmações do artigo: o tom do sopro de Still fica em 120~Hz, e o sopro do VSD tem banda larga. A figura entrega 1333 colunas por segundo, o que limita o conteúdo recuperável a cerca de 667~Hz.

\medskip
\noindent\textbf{Palavras-chave:} processamento digital de imagens; correlação espacial; extração de curvas; sons cardíacos; sopro de Still; comunicação interventricular.
\end{abstract}
\clearpage


\tableofcontents
\clearpage

\section{Introdução}

\subsection{Contextualização e tema}

A ausculta é a primeira ferramenta no exame do coração. Em cada batida o ouvido encontra dois sons principais: o primeiro som (\Sone), do fechamento das valvas entre os átrios e os ventrículos, e o segundo (\Stwo), do fechamento das valvas de saída. O intervalo entre eles é a sístole. Um sopro é um som extra, ouvido na sístole ou na diástole. Na sístole, ele pode ser inocente, sem doença cardíaca, ou patológico. O artigo de Tomas, Rončević e Zelenika~\cite{tomas2011} compara três casos em crianças: sem sopro (Normal), com o sopro inocente de Still e com o sopro patológico da comunicação interventricular (VSD, do inglês \emph{ventricular septal defect}). Segundo o artigo, o sopro de Still é um tom musical de banda estreita, com 90 a 170~Hz (média de 118,75~Hz), e o do VSD é um ruído áspero de banda larga, de 80 a 500~Hz.

O artigo mostra os três sinais em sua Figura~1, como gráficos da amplitude em função do tempo. Eles só existem ali como imagem: uma imagem com três gráficos sobrepostos a uma grade, com eixos, números, legendas e rótulos (\Sone, \Stwo, \emph{systole}, SM, VSD). Este trabalho trata esses gráficos como imagens digitais e recupera deles os sinais de áudio. É um problema de processamento digital de imagens, porque o sinal está codificado na posição vertical de uma curva a cada coluna de pixels, e é preciso separar a curva de tudo o que a cerca.

\subsection{Fundamentação teórica}

\paragraph{Imagem digital e cor.}
Uma imagem colorida em RGB é uma matriz de pixels com três canais (R, G e B), cada um com valores de 0 a 255 em 8 bits. A Figura~1 do artigo foi ampliada em 200\% e comprimida com perdas, então as bordas dos traços são suavizadas (antialiasing) e há pequenas variações de cor sobre os fundos planos. Um pixel de borda é uma mistura da cor do traço com a cor do fundo: $p = a\,c + (1-a)\,f$, em que $c$ é a cor do traço, $f$ a do fundo e $a \in [0,1]$ a cobertura do pixel pelo traço.

\paragraph{Correlação e convolução bidimensionais.}
A filtragem espacial calcula cada pixel de saída como uma soma ponderada da vizinhança do pixel de entrada. Para uma imagem $I$ e uma máscara $h$ de $M\times N$ elementos, a correlação não normalizada com preenchimento por zeros é
\begin{equation}
g(i,j) \;=\; \sum_{m=0}^{M-1}\sum_{n=0}^{N-1} h(m,n)\; I\bigl(i+m-a_y,\; j+n-a_x\bigr),
\qquad a_y = \lfloor M/2\rfloor,\; a_x = \lfloor N/2\rfloor,
\label{eq:correlacao}
\end{equation}
em que $I(y,x)=0$ fora da imagem. A saída tem o mesmo tamanho da entrada. A convolução é a mesma soma com a máscara girada de $180^\circ$, então as duas operações coincidem quando a máscara tem simetria central, como a média e o filtro gaussiano. Para a máscara de Sobel, que é antissimétrica, a diferença é só o sinal da resposta~\cite{gonzalez2018}.

\paragraph{Filtro gaussiano e filtro de Sobel.}
O filtro gaussiano suaviza a imagem e reduz o ruído. A máscara é a amostragem de
\begin{equation}
G(x,y) = \frac{1}{2\pi\sigma^2}\, e^{-(x^2+y^2)/(2\sigma^2)},
\end{equation}
normalizada para somar 1, o que preserva o brilho médio, com lado $2\lceil 3\sigma\rceil+1$, que cobre $\pm 3\sigma$. Ela é separável, pois $G(x,y)=G_1(x)\,G_1(y)$. O filtro de Sobel estima as derivadas parciais da imagem com duas máscaras $3\times 3$ que somam zero, de modo que a resposta é nula em regiões de brilho constante:
\begin{equation}
S_x=\begin{bmatrix}-1&0&1\\-2&0&2\\-1&0&1\end{bmatrix},\qquad
S_y=S_x^{\mathsf T}=\begin{bmatrix}-1&-2&-1\\0&0&0\\1&2&1\end{bmatrix}.
\end{equation}
Com $g_x$ e $g_y$ as respostas, a magnitude do gradiente é $\sqrt{g_x^2+g_y^2}$ e a direção é $\operatorname{atan2}(g_y,g_x)$.

\paragraph{Limiarização, componentes conexos e afinamento.}
O método de Otsu~\cite{otsu1979} escolhe o limiar $t$ que maximiza a variância entre as duas classes de pixels que ele separa, $\sigma_B^2(t)=\omega_0(t)\,\omega_1(t)\,[\mu_0(t)-\mu_1(t)]^2$, em que $\omega_k$ e $\mu_k$ são a proporção e a média de cada classe. Os componentes conexos de uma imagem binária são os grupos de pixels ligados entre si por vizinhança de 8 pixels. O afinamento (esqueletização) reduz um traço grosso a uma linha de 1 pixel; usamos o algoritmo de Lee~\cite{lee1994}.

\paragraph{Amostragem, quantização e espectro.}
Um sinal contínuo amostrado a $f_s$ amostras por segundo só representa componentes abaixo de $f_s/2$ (frequência de Nyquist)~\cite{oppenheim2010}. Na figura, a amostragem é espacial: cada coluna de pixels é uma amostra do tempo, e a resolução horizontal da imagem limita a banda do sinal recuperável. A quantização em 16 bits (PCM com sinal) leva uma amostra $x\in[-1,1]$ ao inteiro $\operatorname{round}(32767\,x)$. O espectro de um trecho é estimado pelo periodograma com janela de Hann, $P(f)\propto|\mathrm{DFT}(x\cdot w)|^2$, e a evolução do espectro no tempo, pelo espectrograma, que é o mesmo cálculo em janelas curtas e deslizantes.

\subsection{Objetivos}

O objetivo geral é automatizar a extração e a reconstrução dos sinais de áudio dos três traçados da Figura~1 do artigo~\cite{tomas2011}, aplicando filtragem espacial, limiarização, remoção de elementos gráficos e amostragem de curvas. Os objetivos específicos são:
\begin{enumerate}
\item implementar uma função de correlação bidimensional não normalizada, com preenchimento por zeros, no domínio do espaço, sem funções prontas de correlação ou filtragem, e testá-la com as imagens fornecidas (Requisito~1);
\item isolar automaticamente cada um dos três gráficos (Normal, Still e VSD);
\item remover a grade, os eixos, as marcas de escala, os textos e as anotações;
\item reduzir o ruído e realçar a estrutura da curva com a correlação do Requisito~1, com máscaras gaussianas e de Sobel;
\item mapear os pixels da curva para o tempo e a amplitude, reamostrar a 5000 amostras por segundo e quantizar em 16 bits, em um único canal;
\item gravar \texttt{normal.wav}, \texttt{still.wav} e \texttt{vsd.wav}, e comparar os três sinais na forma de onda, no espectro e na audição (Requisito~2).
\end{enumerate}

\section{Materiais e métodos}

\subsection{Dados de entrada}

A entrada principal é a \texttt{Figure1.png}, a Figura~1 do artigo~\cite{tomas2011} ampliada em 200\%: uma imagem RGB de $707\times 867$ pixels, 8 bits por canal, com 5837 cores distintas (a compressão com perdas e o antialiasing geram muitos tons intermediários) e 73\% de fundo branco. Ela traz três gráficos empilhados, rotulados (a) Normal, (b) Still e (c) VSD. Cada gráfico tem uma grade cinza, eixos com números, uma curva azul, anotações em azul (\Sone, \Stwo, \emph{systole}, SM e VSD) e a legenda preta ``Diaphragm mode; Normalized''. Os rótulos de tempo de cada gráfico vão de 0,05 em 0,05~s, em 10 intervalos, e cobrem 0,5~s: 5,665 a 6,165~s em (a), 1,345 a 1,845~s em (b) e 1,085 a 1,585~s em (c). Esses valores foram lidos nos rótulos da figura e ficam em \texttt{src/config.py}, sem reconhecimento óptico de caracteres. O eixo vertical vai de $-75\%$ a $+75\%$ em passos de 25\%.

Para o Requisito~1, as imagens fornecidas são a \texttt{Shapes.png} ($800\times 400$, só preto e branco, com um triângulo e um quadrado com um losango) e a \texttt{testpat.1k.color2.tif} ($1024\times 1024$, colorida, com 346 cores e uma foto no centro).

Antes de escrever o processamento, medi a estrutura da figura com um programa de exploração. As premissas do planejamento que não se confirmaram foram:
\begin{itemize}
\item a escala vertical muda de um gráfico para outro: 1,14 pixel por 1\% em (a) e cerca de 1,06 em (b) e (c), e o gráfico (c) fica cerca de 3,6 pixels mais à esquerda, então a calibração precisa ser feita por gráfico;
\item a linha de 0\% quase não aparece, porque a curva a cobre, e precisa ser deduzida das linhas de $\pm25$, $\pm50$ e $\pm75\%$;
\item a curva e as anotações têm a mesma cor azul, mas as anotações são componentes separados da curva (a menor distância é de 6,3 pixels);
\item nos gráficos (a) e (c) há colunas com mais de 100 pixels azuis, nos picos verticais, e isso invalida o uso do centroide da coluna como regra de extração;
\item a figura só tem cerca de 1333 colunas por segundo, e o artigo grava a 8000~Hz.
\end{itemize}

\subsection{Ferramentas e conhecimentos utilizados}

O código é escrito em Python 3.13, com NumPy 2.5.4 para o cálculo numérico, OpenCV 5.0 para ler e gravar imagens, rotular componentes conexos e desenhar retângulos, scikit-image 0.26 para o limiar de Otsu e o afinamento~\cite{otsu1979,lee1994}, SciPy 1.18 para a opção de interpolação cúbica, e Matplotlib 3.11 para as figuras. Os testes usam o pytest. O relatório é composto em \LaTeX. O trabalho foi planejado e dividido em 14 tarefas no GitHub (uma \emph{issue} cada, em um \emph{milestone}), cada uma entregue em uma \emph{branch} própria, com \emph{pull request}.

Foram escritos do zero: a correlação bidimensional, as máscaras (média, gaussiana, Sobel), o cálculo do gradiente e a filtragem de suavização do Requisito~1; a segmentação dos gráficos, a remoção da grade e das anotações, a extração e a calibração da curva; a reamostragem, a quantização e a escrita do \texttt{.wav} (pelo módulo \texttt{wave} da biblioteca padrão); e as métricas de nível, espectro e espectrograma, que usam só \texttt{numpy.fft}. As funções \texttt{scipy.signal.correlate2d} e \texttt{cv2.filter2D} aparecem somente nos testes, como referência de validação, e um teste automático varre o código e falha se elas, ou filtros prontos equivalentes, forem usadas no processamento.

\subsection{Correlação bidimensional (Requisito 1)}

A função \texttt{correlate2d\_manual(I, h)} implementa a Equação~\ref{eq:correlacao}. Ela aceita uma imagem 2D de qualquer tipo numérico e uma máscara de qualquer tamanho, par ou ímpar, converte tudo para ponto flutuante de 64 bits e devolve uma matriz do mesmo tamanho de $I$. A âncora fica em $(\lfloor M/2\rfloor,\lfloor N/2\rfloor)$, o mesmo critério do \texttt{cv2.filter2D}. Para uma máscara par, o preenchimento de zeros é de $\lfloor M/2\rfloor$ linhas acima e $M-1-\lfloor M/2\rfloor$ abaixo, e o análogo nas colunas.

Foram escritas duas versões. A primeira, \texttt{correlate2d\_naive}, segue a definição com quatro laços (pixel de saída e elemento da máscara) e serve de referência. A segunda é a usada no pipeline, e troca a ordem dos laços: a imagem é copiada para uma matriz com bordas de zeros, e a saída é a soma, sobre os $M N$ elementos da máscara, da imagem deslocada multiplicada pelo coeficiente:
\begin{center}
\begin{minipage}{0.8\textwidth}
\small
\begin{verbatim}
padded = zeros(H + acima + abaixo, W + esq + dir)
padded[acima:acima+H, esq:esq+W] = I
saida  = zeros(H, W)
para m em 0..M-1:
    para n em 0..N-1:
        saida += h[m, n] * padded[m : m+H, n : n+W]
\end{verbatim}
\end{minipage}
\end{center}
Os dois laços externos percorrem a máscara e a soma interna é vetorizada, o que mantém a operação no domínio do espaço e o custo em $O(H\,W\,M\,N)$ (como na definição), mas a menos de 0,1~s para cinco filtros em uma imagem de $1024\times 1024$. A função recusa entradas que não sejam 2D ou que estejam vazias.

\subsection{Máscaras de filtragem}

O módulo \texttt{kernels} gera as máscaras usadas: a de média $n\times n$ com todos os elementos iguais a $1/n^2$; a gaussiana, normalizada para somar 1, com lado $2\lceil 3\sigma\rceil+1$ (7$\times$7 para $\sigma=1$); e as de Sobel $S_x$ e $S_y$. Com elas, \texttt{suavizar\_gaussiano} e \texttt{gradiente\_sobel} dão a imagem suavizada, as componentes $g_x$ e $g_y$, a magnitude e a direção do gradiente.

\subsection{Pipeline de reconstrução}

A Figura~\ref{fig:fluxo} resume o fluxo. O programa \texttt{main.py} executa tudo em sequência, grava os arquivos \texttt{.wav} e as imagens de cada etapa, e todos os parâmetros ajustáveis ficam em \texttt{src/config.py}.

\begin{figure}[H]
\centering
\begin{tikzpicture}[node distance=5mm, font=\footnotesize,
  caixa/.style={draw, rounded corners=2pt, align=center, minimum height=9mm, minimum width=22mm, fill=gray!8},
  seta/.style={->, thick}]
\node[caixa] (a) {Figure1.png};
\node[caixa, right=of a] (b) {1.~Segmentar\\os gráficos};
\node[caixa, right=of b] (c) {2.~Remover\\a grade};
\node[caixa, right=of c] (d) {3.~Remover\\anotações};
\node[caixa, below=7mm of a] (e) {4.~Suavizar\\e binarizar};
\node[caixa, right=of e] (f) {5.~Extrair\\a curva};
\node[caixa, right=of f] (g) {6.~Calibrar\\tempo e amplitude};
\node[caixa, right=of g] (h) {7.~Reamostrar\\e gravar .wav};
\draw[seta] (a)--(b); \draw[seta] (b)--(c); \draw[seta] (c)--(d);
\draw[seta] (d.south) -- ++(0,-3.5mm) -| ($(e.north)+(0,3.5mm)$) -- (e.north);
\draw[seta] (e)--(f); \draw[seta] (f)--(g); \draw[seta] (g)--(h);
\end{tikzpicture}
\caption{Etapas do pipeline. A suavização da etapa~4 usa a correlação do Requisito~1.}
\label{fig:fluxo}
\end{figure}

\paragraph{1. Segmentação dos gráficos.}
Os gráficos são localizados sem coordenadas fixas. Cada um tem à esquerda uma barra cinza com os rótulos de amplitude, que as legendas ``(a)'', ``(b)'', ``(c)'' e a descrição da figura não têm. A barra é a primeira faixa de colunas, com pelo menos 10 pixels de largura, em que mais da metade dos pixels não é branca. Os gráficos são as faixas de linhas, com pelo menos 10\% da altura da imagem, em que a barra é preenchida. Dentro de cada faixa, a faixa dos rótulos de tempo é o último trecho de pelo menos 8 linhas em que mais de 35\% dos pixels são cinza. As linhas horizontais da grade são as linhas com mais de 70\% de pixels cinza, e as verticais são as colunas com mais de metade da altura em cinza, com trechos separados por até 3 pixels unidos. A posição das 10 linhas verticais em função do índice, ajustada por uma reta, dá as bordas do eixo de tempo, $x_{\text{esq}}$ e $x_{\text{dir}}$. As linhas horizontais, indexadas a partir da de $-75\%$ (a última), dão por ajuste a posição das linhas de $+75\%$, $0\%$ e $-75\%$ ($y_{+75}$, $y_0$ e $y_{-75}$), inclusive a de 0\%, que a curva cobre. A região de interesse (ROI) vai de $x_{\text{esq}}+2$ a $x_{\text{dir}}-1$ e, na vertical, da moldura superior (um passo de 25\% acima da linha de $+75\%$) até a faixa de rótulos de tempo. Se a detecção falhar, o programa usa as coordenadas medidas, registra um aviso e continua.

\paragraph{2. Remoção da grade, dos eixos e dos textos pretos.}
A grade, os eixos, os números e a legenda são neutros ($R=G=B$), e a curva e as anotações são azuis. Define-se o ``azulado'' de um pixel como $\beta=\max\bigl(0,\;B-\max(R,G)\bigr)$. Para um pixel que mistura a curva ($c$) com um fundo neutro ($f$, de valor $\varphi$ nos três canais), $p=a\,c+(1-a)\,f$, e portanto
\begin{equation}
\beta(p)=a\,\bigl[c_B-\max(c_R,c_G)\bigr]=a\,\beta(c),
\end{equation}
pois o termo $(1-a)\varphi$ aparece igual em $B$ e em $\max(R,G)$ e se cancela. O azulado de um pixel da borda da curva é proporcional à cobertura $a$, e não depende de ela estar sobre o branco ou sobre uma linha da grade. Nos pixels neutros da figura, o ruído de compressão faz $\beta$ chegar a 9, então o limiar é 12. A cobertura da curva é $\operatorname{clip}\bigl((\beta-12)/(\beta_{99}-12),0,1\bigr)$, em que $\beta_{99}$ é o percentil 99 de $\beta$ entre os pixels com $\beta>12$. A imagem resultante é a cobertura desenhada em tinta escura sobre branco.

\paragraph{3. Remoção das anotações.}
Os componentes conexos (vizinhança de 8) da máscara de cobertura são classificados pela forma. A curva é o de maior área. Os outros são anotações se tiverem no máximo 20 pixels de altura e 45 de largura, e se ocuparem pelo menos 50\% da sua caixa envolvente (as letras são compactas, e um pedaço do traço da curva é fino e longo). As anotações são apagadas e as suas caixas guardadas, porque a análise usa as caixas de \Sone{} e \Stwo{} (Seção~\ref{sec:metodo-analise}). Se a curva restante não cobrir 98\% da largura da ROI, o programa registra um aviso.

\paragraph{4. Suavização e binarização.}
A cobertura da curva é suavizada com a máscara gaussiana ($\sigma=1$, 7$\times$7) pela função de correlação do Requisito~1, e binarizada pelo limiar de Otsu. O mapa de bordas de Sobel é calculado sobre a imagem suavizada e o esqueleto de 1 pixel (algoritmo de Lee~\cite{lee1994}) é calculado sobre a máscara binária. Ambos servem para visualizar as etapas do processamento. O valor de $\sigma$ foi escolhido pela comparação de cinco valores (Seção~\ref{sec:res-realce}).

\paragraph{5. Extração e calibração da curva.}
Em cada coluna da máscara, o topo e a base da curva são as linhas do primeiro e do último pixel. Colunas sem curva seriam interpoladas, mas nenhuma ocorreu. O tempo e a amplitude vêm das medidas da etapa~1:
\begin{equation}
t(x)=t_0+\bigl(x-x_{\text{esq}}\bigr)\,\frac{0{,}5}{x_{\text{dir}}-x_{\text{esq}}},\qquad
A(y)=\bigl(y_0-y\bigr)\,\frac{150}{y_{-75}-y_{+75}}\;[\%],
\end{equation}
em que $t_0$ é o tempo do rótulo da borda esquerda do gráfico. Cada gráfico tem a sua calibração.

Em cada coluna, é preciso escolher um valor de amplitude entre o topo e a base. A regra escolhida (``picos'') usa o ponto médio, que acompanha bem as subidas e descidas, mas troca o ponto médio pelo extremo nas colunas de virada. Só as colunas cuja extensão $\text{topo}-\text{base}$ passa de 6 pixels (o traço mede de 3 a 4) podem ser trocadas. Se o topo é um máximo local entre as colunas vizinhas e a subida do topo sobre os vizinhos é pelo menos igual à descida da base abaixo deles, vale o topo. Senão, se a base é um mínimo local, vale a base. O ponto médio sozinho (regra ``centro'') perde cerca de um terço da amplitude dos picos, como mostram os testes sintéticos da Seção~\ref{sec:res-extracao}.

\paragraph{6. Reamostragem e gravação.}
O sinal extraído tem uma amostra por coluna (cerca de 664 por gráfico). Ele é interpolado linearmente na grade uniforme de $f_s=5000$~Hz, $t_k=t_0+k/f_s$ com $k=0,\dots,2499$, que dá 0,5~s. As pontas, que a ROI deixa de fora, repetem o valor da borda. Depois, retira-se o nível médio (DC), o sinal é normalizado de modo que o maior valor absoluto seja 0,95 do fundo de escala, e quantizado por $\operatorname{round}(32767\,x)$ em 16 bits. O arquivo é um WAV PCM mono. A interpolação cúbica ficou como opção, porque ultrapassa os picos em cerca de 10\%.

\subsection{Validação}

Como não há áudio original para comparar, a extração foi validada pela ida e volta: sinais sintéticos conhecidos, a 8000~Hz como no artigo, foram desenhados como na figura (traço de 2,5 pixels, antialiasing, 1333 colunas por segundo e 1,06 pixel por 1\%), passaram pelas mesmas etapas de suavização, binarização, extração e reamostragem, e foram comparados ao sinal original a 5000~Hz pela correlação, pelo erro quadrático médio relativo (NRMSE) e pela razão entre as amplitudes dos picos. Há três tipos: dois estouros amortecidos (\emph{burst}, como \Sone{} e \Stwo), os estouros com ruído de 100 a 350~Hz entre eles (\emph{sopro}), e os estouros com ruído de 150 a 600~Hz (\emph{denso}, como o gráfico (c)). A suíte de testes automáticos, com mais de 300 testes, cobre cada módulo com casos conhecidos, compara a correlação com as funções prontas nos testes, e confere o resultado em cada uma das três figuras reais.

\subsection{Análise dos sinais reconstruídos}
\label{sec:metodo-analise}

Cada sinal (0,5~s, sem DC, em \% da escala do gráfico) é dividido em janelas: \Sone{} e \Stwo{} (80~ms cada, centradas no evento), o sopro sistólico (entre elas) e a diástole (antes de \Sone{} e depois de \Stwo, usada como nível de fundo). Os eventos \Sone{} e \Stwo{} são localizados pelos rótulos da própria figura: a caixa de anotação mais à esquerda é \Sone{}, a mais à direita é \Stwo, e o evento é o máximo da envolvente RMS (janela de 10~ms) de 10~ms antes a 50~ms depois do centro do rótulo. Em cada janela, medem-se o pico, o valor RMS e o fator de crista (pico sobre RMS), e as razões em decibéis $20\log_{10}(\text{RMS}_1/\text{RMS}_2)$ entre o sopro, a diástole e \Sone. No sopro, mede-se também a periodicidade: o maior pico da autocorrelação normalizada para períodos de 2,5 a 16,7~ms (60 a 400~Hz).

O espectro de cada janela é o periodograma com janela de Hann e preenchimento por zeros até 8192 pontos, entre 40 e 660~Hz. Dele saem a frequência dominante, o centroide, a largura a meia potência do pico dominante, a fração da energia em quatro faixas (40--80, 80--200, 200--500 e 500--660~Hz) e a planura espectral, que é a razão entre as médias geométrica e aritmética da potência de 80 a 500~Hz (perto de 0 para um tom puro, e perto de 0,56 para o ruído branco). O espectrograma usa janelas de Hann de 25~ms com passo de 2~ms.

\section{Resultados}
% seção em redação

\section{Discussão}

\subsection{Problemas e dificuldades encontrados}

\paragraph{A resolução da imagem limita o áudio.}
A figura tem cerca de 1333 colunas por segundo, e uma coluna é uma amostra. Pelo critério de Nyquist, só dá para recuperar o que está abaixo de uns 667~Hz, ao passo que o artigo gravou a 8000~Hz. A especificação pede 5000~Hz, e a reamostragem para essa taxa só interpola: ela não cria detalhe, e os arquivos de 5~kHz contêm, na prática, o conteúdo de 0 a 667~Hz. O sopro de Still (120~Hz) e o do VSD (80 a 500~Hz, segundo o artigo) cabem nessa banda, e por isso a comparação ainda faz sentido. O que estiver acima de 667~Hz, em \Sone, em \Stwo{} ou em qualquer outro trecho, simplesmente não está na figura.

\paragraph{O gráfico denso.}
O gráfico (c) tem oscilações que se sobrepõem na mesma coluna. A extração devolve uma amostra por coluna, então só o envelope dessas oscilações é recuperado. No teste sintético equivalente, a correlação com o sinal original é de 0,78 e o erro relativo é de 0,67, contra 0,98 e 0,19 nos estouros isolados. Isso quer dizer que o \texttt{vsd.wav} é uma aproximação mais pobre do que o \texttt{normal.wav} e o \texttt{still.wav}. Em um experimento exploratório, fora do código final, uma extração com mais de uma amostra por coluna (o topo e a base em ordem) manteve o pico, mas acrescentou ruído de alta frequência e piorou a correlação (0,73, 0,69 e 0,46 nos três tipos de sinal sintético), e foi descartada.

\paragraph{O centroide e o esqueleto não servem para extrair a amplitude.}
O centroide da coluna erra nos picos verticais (em alguns trechos, a coluna tem mais de 100 pixels de curva), e o ponto médio entre o topo e a base perde cerca de um terço do pico. Foi preciso uma regra que preserve o extremo nas colunas de virada. O esqueleto de 1 pixel, que seria o resultado natural do afinamento pedido, erra os extremos em 3 a 8 pixels, e por isso ficou só para visualização.

\paragraph{A linha de 0\% está coberta pela curva.}
A grade horizontal mostra as linhas de $\pm25$, $\pm50$ e $\pm75\%$, mas a curva cobre a de 0\%. A posição dela vem do ajuste das linhas vizinhas, com a hipótese de que são igualmente espaçadas. Foi confirmada na única linha residual que aparece no gráfico (a), a 1 pixel do valor deduzido, mas é uma dedução, e não uma medida.

\paragraph{A cor não separa a curva das anotações.}
A curva e as anotações têm o mesmo azul. A remoção por componentes conexos funciona porque nenhuma anotação toca a curva (a mais próxima fica a 6,3 pixels). Se uma tocasse, as duas virariam um só componente, e o programa só avisaria que a curva não cobre a largura inteira, sem separá-las. A suavização gaussiana, que poderia fundir uma anotação à curva, é feita depois da remoção.

\paragraph{Os níveis absolutos não são comparáveis.}
A legenda diz que os gráficos estão normalizados, mas os picos de 93\%, 55\% e 88\% mostram que não é pelo máximo de cada trecho, e o critério não é conhecido. Só valem as razões dentro de cada sinal (por exemplo, o sopro sobre \Sone). O volume relativo entre os arquivos de áudio também se perde com a normalização por pico: o programa devolve os níveis originais para a análise, e oferece a normalização por escala, que o preserva.

\paragraph{Não há áudio original para comparar.}
O resultado não pôde ser comparado com a gravação que deu origem à figura. A validação usou sinais sintéticos de gabarito conhecido, passados por todo o pipeline, e a conferência com as afirmações do artigo. Os sinais sintéticos são mais simples do que os reais, então os erros medidos nesse teste são uma estimativa.

\subsection{Comentários críticos sobre os resultados}

Os resultados são consistentes entre si em vários níveis. A curva extraída coincide com o traço da figura em cada gráfico, os picos de \Sone{} e \Stwo{} do gráfico (a) batem com a leitura direta (cerca de 75\% e de 92\% a $-72\%$), o tom do sopro de Still em 120~Hz coincide com a média do artigo (118,75~Hz), e o sopro do VSD tem a banda larga que o artigo descreve, com 46\% da energia entre 200 e 500~Hz. Sete das oito afirmações conferidas se reproduzem. A que não se reproduz (o pico do VSD acima de 200~Hz) é uma estatística de 20 gravações, e aqui há uma só.

Esses resultados têm limites que convém dizer. A análise cobre uma gravação de cada tipo, então ela mostra que os sinais reconstruídos são coerentes com o artigo, e não que o método classifica sopros. A localização de \Sone{} e \Stwo{} usa a posição dos rótulos da própria figura, e foi conferida contra a leitura visual dentro de 20~ms; se o layout fosse outro, a regra teria de ser revista. A audição foi informal e feita por uma só pessoa. O Normal e o VSD soaram como as medidas previam, mas o Still soou diferente do esperado: em vez de um zumbido suave entre as duas batidas, o ouvinte descreveu batidas a mais. As medidas são compatíveis com isso, pois o sopro de Still é periódico (repetição a cada 8,4~ms) e só 7,5~dB mais fraco que \Sone, mas essa explicação não foi testada com ouvintes e sons de controle.

Por fim, uma observação sobre o método. A etapa mais importante para a qualidade do áudio não foi a filtragem espacial pedida pelo enunciado (a correlação com o gaussiano e o Sobel), e sim a regra de extração por coluna: a mesma máscara, extraída com o ponto médio, perde um terço da amplitude dos picos. O gaussiano e o Sobel cumprem o papel de limpar e realçar a curva, mas o $\sigma$ tem de ser pequeno (1,0), porque valores maiores alongam os picos em 5 a 18 pixels e distorcem o sinal.

\subsection{Possíveis melhorias}

Com uma imagem de maior resolução, ou com o sinal original, a banda recuperável seria maior. No gráfico denso, uma extração que acompanhe o traço (seguindo a curva de pixel em pixel, em vez de reduzir cada coluna a um valor) poderia recuperar oscilações que hoje se perdem. Para a análise, um teste de audição com vários ouvintes e sons de controle diria se a pulsação do sopro de Still é de fato o que se ouve como batidas extras.

\section{Conclusão}
% seção em redação


\begin{thebibliography}{9}
\bibitem{tomas2011} B.~Tomas, Ž.~Rončević e D.~Zelenika. Classification of Pathologic and Innocent Heart Murmur Based on Multimedia Presentations of Acoustic Heart Signals. In: \emph{CONTENT 2011: The Third International Conference on Creative Content Technologies}. IARIA, 2011. ISBN 978-1-61208-157-1.
\bibitem{gonzalez2018} R.~C. Gonzalez e R.~E. Woods. \emph{Digital Image Processing}. 4.\,ed. Pearson, 2018.
\bibitem{otsu1979} N.~Otsu. A threshold selection method from gray-level histograms. \emph{IEEE Transactions on Systems, Man, and Cybernetics}, 9(1):62--66, 1979.
\bibitem{lee1994} T.~C. Lee, R.~L. Kashyap e C.~N. Chu. Building skeleton models via 3-D medial surface/axis thinning algorithms. \emph{CVGIP: Graphical Models and Image Processing}, 56(6):462--478, 1994.
\bibitem{oppenheim2010} A.~V. Oppenheim e R.~W. Schafer. \emph{Discrete-Time Signal Processing}. 3.\,ed. Pearson, 2010.

\end{thebibliography}

\end{document}
